% Source: https://artofproblemsolving.com/community/c3h2538457_polynominal_proof
% Stable identity: aops-topic-2538457-q01
\textbf{(a)} Since $c=rs\ne0$, both roots are nonzero, and $\dfrac1r+\dfrac1s=\dfrac{s+r}{rs}$. By Vieta's formulas $r+s=-b$ and $rs=c$, so the common value is $\boxed{-\dfrac bc}$. $\blacksquare$

\textbf{(b)} For the roots $1,2,3$: $\dfrac{1+2+3}{1\cdot2\cdot3}=1$, but $\dfrac11+\dfrac12+\dfrac13=\dfrac{11}6\ne1$. So the claim fails.

In general, over a common denominator,
\[\dfrac1r+\dfrac1s+\dfrac1t=\dfrac{st+tr+rs}{rst}.\]
The numerator is the sum of the products of the roots \emph{two at a time}, not their sum; for two roots the two coincide, which is why (a) works. By Vieta's formulas for $x^3+bx^2+cx+d$, $rs+st+tr=c$ and $rst=-d$, so
\[\dfrac1r+\dfrac1s+\dfrac1t=\boxed{-\dfrac cd}.\]
(Check: for $x^3-6x^2+11x-6$ this gives $\dfrac{11}6$.)
